\begin{frame}{과적합을 숫자로: 10단계 vs 500단계}
  \begin{columns}[T]
    \begin{column}{0.52\textwidth}
      \centering
      \includegraphics[width=\linewidth]{ch07_gbdt_overfitting.png}
      \vskip0.3em
      {\footnotesize train MSE는 계속 0.005까지, test MSE는 10단계 최저 후
      1.290까지 다시 올라감.}
    \end{column}
    \begin{column}{0.45\textwidth}
      세 숫자가 전부 핵심:
      \begin{enumerate}
        \item \textbf{train MSE는 0.005까지 계속 하락} — 500개 트리가
          학습 데이터의 잡음 하나까지 \textbf{거의 완벽히 외움}
        \item \textbf{test MSE는 10단계에서 최저(0.727)}
        \item \textbf{500단계까지 하면 1.290까지 다시 상승} — 최저점 대비
          \textbf{약 두 배}로 악화, ``더 오래 학습한 모델''이 정점의
          \textbf{절반 수준}
      \end{enumerate}
    \end{column}
  \end{columns}
  \vskip0.5em
  \textbf{왜 GBDT만?} RF의 101번째 트리는 데이터를 \textbf{처음부터 다시}
  보고 다른 트리의 노이즈 오류를 ``따라잡아야 할 의무''가 없다(다수결이
  희석). GBDT의 101번째 트리는 앞 100개 합이 \textbf{아직 못 맞춘 것
  (= 학습 데이터의 노이즈 잔차)을 \textbf{의도적으로} 목표}로 삼는다 —
  적당히 쌓이면 잔차 대부분은 \textbf{라벨링 노이즈}인데, GBDT는 그것을
  ``정당한 학습 목표''로 받아들인다.
\end{frame}
